\appendix

% ══════════════════════════════════════════════════════════════════════════════
\chapter{Full HIGGS Scaling Experiment Results}
\label{app:scaling}

\Cref{tab:app_scaling_wf1,tab:app_scaling_co2} provide the complete numeric results for the HIGGS subset scaling experiment introduced in \Cref{design_scaling} and discussed in \Cref{sec:results_scaling}. \Cref{fig:scaling_higgs} in the main text shows these values as line charts; the tables below give exact figures for reproducibility. Random Forest is excluded from subset sizes above 500{,}000 due to prohibitive training times. All CO\textsubscript{2} values are hardware-corrected (see \Cref{implementation_emissions}).

\begin{table}[htbp]
\centering
\caption{Weighted F1-score per model and HIGGS subset size. Values are means across three random seeds. Dashes indicate runs that were not performed.}
\label{tab:app_scaling_wf1}
\begin{tabular}{@{} r rrrrr @{}}
\toprule
\textbf{Rows} & \textbf{Log.\ Reg.} & \textbf{Random Forest} & \textbf{XGBoost (CPU)} & \textbf{XGBoost (GPU)} & \textbf{MLP} \\
\midrule
1{,}000       & 0.630 & 0.660 & 0.663 & 0.653 & 0.599 \\
10{,}000      & 0.628 & 0.711 & 0.701 & 0.704 & 0.653 \\
50{,}000      & 0.633 & 0.715 & 0.713 & 0.712 & 0.697 \\
100{,}000     & 0.636 & 0.723 & 0.719 & 0.717 & 0.718 \\
200{,}000     & 0.635 & 0.726 & 0.726 & 0.726 & 0.732 \\
500{,}000     & 0.638 & 0.733 & 0.738 & 0.738 & 0.751 \\
1{,}000{,}000 & 0.638 & ---   & 0.745 & 0.744 & 0.760 \\
5{,}000{,}000 & 0.637 & ---   & 0.757 & 0.757 & 0.778 \\
11{,}000{,}000 & 0.637 & ---  & 0.760 & 0.760 & 0.785 \\
\bottomrule
\end{tabular}
\end{table}

\begin{sidewaystable}[htbp]
\centering
\caption{Hardware-corrected CO\textsubscript{2} emissions (mg) and training time (s) per model and HIGGS subset size. Values are means across three random seeds; the \gls{mlp} at 5M and 11M rows was run under a single seed. The \gls{mlp} at 11M rows (73{,}145\,mg\,=\,73.1\,g) is the study's highest single training cost.}
\label{tab:app_scaling_co2}
\setlength{\tabcolsep}{4pt}
\begin{tabular}{@{} r rr rr rr rr rr @{}}
\toprule
 & \multicolumn{2}{c}{\textbf{Log.\ Reg.}} & \multicolumn{2}{c}{\textbf{Random Forest}} & \multicolumn{2}{c}{\textbf{XGBoost (CPU)}} & \multicolumn{2}{c}{\textbf{XGBoost (GPU)}} & \multicolumn{2}{c}{\textbf{MLP}} \\
\cmidrule(lr){2-3}\cmidrule(lr){4-5}\cmidrule(lr){6-7}\cmidrule(lr){8-9}\cmidrule(lr){10-11}
\textbf{Rows} & CO\textsubscript{2} (mg) & Time (s) & CO\textsubscript{2} (mg) & Time (s) & CO\textsubscript{2} (mg) & Time (s) & CO\textsubscript{2} (mg) & Time (s) & CO\textsubscript{2} (mg) & Time (s) \\
\midrule
1{,}000        & 12   & 0.0    & 18    & 0.5    & 16   & 0.2    & 24   & 0.7    & 20   & 0.8    \\
10{,}000       & 12   & 0.0    & 29    & 1.0    & 22   & 0.5    & 33   & 1.2    & 30   & 1.3    \\
50{,}000       & 16   & 0.2    & 105   & 5.4    & 39   & 1.4    & 46   & 1.9    & 100  & 5.7    \\
100{,}000      & 19   & 0.4    & 217   & 11.9   & 44   & 1.8    & 47   & 2.0    & 207  & 12.4   \\
200{,}000      & 22   & 0.9    & 492   & 28.2   & 66   & 3.0    & 51   & 2.2    & 328  & 22.5   \\
500{,}000      & 47   & 3.0    & 1{,}708 & 102.8 & 123 & 6.5    & 64   & 2.7    & 910  & 67.2   \\
1{,}000{,}000  & 92   & 7.0    & ---  & ---    & 226  & 12.7   & 87   & 3.6    & 2{,}038 & 154.7 \\
5{,}000{,}000  & 441  & 39.3   & ---  & ---    & 1{,}243 & 77.4 & 220  & 9.2    & 15{,}588 & 1{,}629 \\
11{,}000{,}000 & 1{,}005 & 92.1 & ---  & ---   & 2{,}703 & 170.6 & 432 & 18.0  & 73{,}145 & 7{,}659 \\
\bottomrule
\end{tabular}
\end{sidewaystable}


% ══════════════════════════════════════════════════════════════════════════════
\chapter{Full MLP Architecture Variation Results}
\label{app:mlp_variation}

\Cref{tab:app_mlp_full} gives the complete numeric results for all twelve \gls{mlp} architecture configurations evaluated in the variation experiment (\Cref{design_mlp_variation,sec:results_mlp_arch}). The heatmaps in \Cref{fig:mlp_variation,fig:mlp_variation_cf1} visualize the same data. All runs use the full 11M-row HIGGS dataset. \gls{cf1} is computed as $\text{CF1} = c_{\mathrm{mg}} / (\text{WF1} \times 100)$, where $c_{\mathrm{mg}}$ is CO\textsubscript{2} in milligrams; lower values indicate a more ecologically efficient configuration.

\begin{table}[htbp]
\centering
\caption{Complete results for the 12 \gls{mlp} architecture configurations on HIGGS (11M rows). Configurations are labelled \textit{depth}$\times$\textit{width}. Training time is single-run wall-clock time. \gls{cf1} is mg CO\textsubscript{2} per WF1 percentage point; lower is better. Bold marks the best value per column.}
\label{tab:app_mlp_full}
\begin{tabular}{@{} l r r r r r @{}}
\toprule
\textbf{Config} & \textbf{WF1} & \textbf{CO\textsubscript{2} (g)} & \textbf{Time (h)} & \textbf{CF1} \\
\midrule
$1\times 128$   & 0.728 & 25.6  & 0.75 & 352 \\
$2\times 128$   & 0.756 & \textbf{13.5} & \textbf{0.39} & \textbf{179} \\
$3\times 128$   & 0.757 & 18.1  & 0.54 & 239 \\
$4\times 128$   & 0.756 & 16.1  & 0.48 & 213 \\
\midrule
$1\times 256$   & 0.736 & 30.3  & 0.95 & 412 \\
$2\times 256$   & 0.768 & 30.5  & 0.93 & 397 \\
$3\times 256$   & 0.771 & 42.6  & 1.26 & 553 \\
$4\times 256$   & 0.772 & 43.5  & 1.28 & 564 \\
\midrule
$1\times 512$   & 0.744 & 48.1  & 1.33 & 647 \\
$2\times 512$   & 0.776 & 54.6  & 1.39 & 704 \\
$3\times 512$   & 0.781 & 80.5  & 2.16 & 1{,}031 \\
$4\times 512$   & \textbf{0.783} & 92.7 & 2.17 & 1{,}184 \\
\bottomrule
\end{tabular}
\end{table}


% ══════════════════════════════════════════════════════════════════════════════
\chapter{Per-Dataset Benchmark Figures}
\label{app:per_dataset}

The three figures below show per-dataset benchmark breakdowns for Wine, Credit, and HIGGS. The combined cross-dataset comparison is presented in \Cref{sec:results_main}; these figures provide a higher-resolution view of each dataset in isolation.

\begin{figure}[htbp]
\centering
\includegraphics[width=0.8\linewidth]{wine_benchmark.pdf}
\caption{Detailed benchmark results for the Wine dataset across all five models.}
\label{fig:app_wine}
\end{figure}

\begin{figure}[htbp]
\centering
\includegraphics[width=0.8\linewidth]{credit_benchmark.pdf}
\caption{Detailed benchmark results for the Credit Card Default dataset across all five models.}
\label{fig:app_credit}
\end{figure}

\begin{figure}[htbp]
\centering
\includegraphics[width=0.8\linewidth]{higgs_benchmark.pdf}
\caption{Detailed benchmark results for the HIGGS dataset. Random Forest is excluded due to prohibitive training times at 11M rows.}
\label{fig:app_higgs}
\end{figure}


% ══════════════════════════════════════════════════════════════════════════════
\chapter{Experimental Setup Reference}
\label{app:setup}

\section{Dataset Characteristics}
\label{app:datasets}

\Cref{tab:app_datasets} summarises the key properties of the three datasets. The feature count excludes the target column; for Credit the identifier column (\texttt{ID}) is additionally dropped before training. Class distribution is measured on the full dataset before any train/test split.

\begin{table}[htbp]
\centering
\caption{Key properties of the three benchmark datasets.}
\label{tab:app_datasets}
\begin{tabular}{@{} l r r r l @{}}
\toprule
\textbf{Dataset} & \textbf{Samples} & \textbf{Features} & \textbf{Classes} & \textbf{Class distribution} \\
\midrule
Wine         &        178 & 13 & 3 & 33.1\,\% / 39.9\,\% / 27.0\,\% \\
Credit       &     30{,}000 & 23 & 2 & 77.9\,\% (no default) / 22.1\,\% (default) \\
HIGGS        & 11{,}000{,}000 & 28 & 2 & 47.0\,\% (background) / 53.0\,\% (signal) \\
\bottomrule
\end{tabular}
\end{table}

\section{Python Software Environment}
\label{app:software}

\Cref{tab:app_software} lists the versions of all core libraries used in this study. A complete \texttt{requirements.txt} is included in the accompanying code repository. All experiments were run on Python 3.13.0 on Windows~11.

\begin{table}[htbp]
\centering
\caption{Core Python library versions used in all experiments.}
\label{tab:app_software}
\begin{tabular}{ll}
\toprule
\textbf{Library} & \textbf{Version} \\
\midrule
Python           & 3.13.0 \\
scikit-learn     & 1.8.0 \\
XGBoost          & 3.2.0 \\
PyTorch          & 2.7.0+cu128 \\
Skorch           & 1.3.1 \\
Optuna           & 4.8.0 \\
CodeCarbon       & 3.2.3 \\
HardwareMonitor  & 1.1.0 \\
pandas           & 3.0.2 \\
NumPy            & 2.4.4 \\
matplotlib       & 3.10.8 \\
seaborn          & 0.13.2 \\
fastparquet      & 2026.3.0 \\
\bottomrule
\end{tabular}
\end{table}


% ══════════════════════════════════════════════════════════════════════════════
\chapter{Supplementary Numeric Results}
\label{app:supplementary}

\section{XGBoost CPU vs.\ GPU}
\label{app:xgb_cpugpu}

\Cref{tab:app_xgb_cpugpu} gives exact values for the XGBoost \gls{cpu} vs.\ \gls{gpu} comparison discussed in \Cref{sec:results_xgb}; \Cref{fig:xgb_cpu_vs_gpu} visualises the same data. Both variants implement the same algorithm and produce equivalent predictive performance; any sub-0.001 difference in WF1 across runs is attributable to floating-point precision differences between \gls{cpu} (float64) and \gls{gpu} (float32) arithmetic and carries no practical meaning. The table therefore omits accuracy and WF1, focusing on the metrics that actually differ: CO\textsubscript{2} emissions, training time, and inference latency. Bold marks the better value within each dataset.

\begin{table}[htbp]
\centering
\caption{Hardware efficiency results for XGBoost \gls{cpu} and \gls{gpu} variants across all three datasets. CO\textsubscript{2} is hardware-corrected training emissions; inference latency is mean per-sample wall-clock time. Predictive performance (WF1) is equivalent across variants on all datasets.}
\label{tab:app_xgb_cpugpu}
\setlength{\tabcolsep}{7pt}
\begin{tabular}{@{} l l r r r r @{}}
\toprule
\textbf{Dataset} & \textbf{Variant} & \textbf{WF1} & \textbf{CO\textsubscript{2} (g)} & \textbf{Time (s)} & \textbf{Latency ($\mu$s)} \\
\midrule
\multirow{2}{*}{Wine}
  & XGBoost (CPU) & 0.972 & \textbf{0.013} & \textbf{0.1}  & \textbf{1464} \\
  & XGBoost (GPU) & 0.972 & 0.017          & 0.3           & 2468 \\
\midrule
\multirow{2}{*}{Credit}
  & XGBoost (CPU) & 0.798 & \textbf{0.021} & \textbf{0.4}  & \textbf{2785} \\
  & XGBoost (GPU) & 0.798 & 0.022          & 0.7           & 4688 \\
\midrule
\multirow{2}{*}{HIGGS}
  & XGBoost (CPU) & 0.760 & 2.763          & 155.2         & \textbf{2021} \\
  & XGBoost (GPU) & 0.760 & \textbf{0.393} & \textbf{16.4} & 3491 \\
\bottomrule
\end{tabular}
\end{table}

\clearpage
\section{Monthly Carbon Intensity Statistics: German Grid 2024}
\label{app:carbon_monthly}

\Cref{tab:app_carbon_monthly} summarises the carbon intensity of the German electricity grid for each calendar month of 2024, derived from 8{,}760 hourly observations obtained via the ElectricityMaps API. These values underpin the seasonal counterfactual analysis in \Cref{sec:results_carbon}. The winter months (October--January) consistently show higher mean intensity than the summer months, driven by reduced solar generation and higher gas and coal dispatch. November records the highest monthly mean; June records the single lowest hourly minimum.

\begin{table}[htbp]
\centering
\caption{Descriptive statistics of hourly carbon intensity (\gls{gco2eq}/\gls{kwh}) for the German bidding zone, January--December 2024. All values are rounded to the nearest integer.}
\label{tab:app_carbon_monthly}
\begin{tabular}{@{} l r r r r r @{}}
\toprule
\textbf{Month} & \textbf{Mean} & \textbf{Median} & \textbf{Min} & \textbf{Max} & \textbf{SD} \\
\midrule
January   & 382 & 380 & 159 & 650 & 127 \\
February  & 344 & 338 & 167 & 576 & 101 \\
March     & 359 & 354 & 143 & 612 & 122 \\
April     & 288 & 272 & 129 & 562 & 101 \\
May       & 290 & 278 & 134 & 509 &  92 \\
June      & 298 & 288 & 119 & 525 & 108 \\
July      & 294 & 289 & 127 & 520 & 107 \\
August    & 320 & 315 & 115 & 536 & 117 \\
September & 322 & 322 & 127 & 546 & 103 \\
October   & 380 & 389 & 147 & 606 & 105 \\
November  & 435 & 456 & 143 & 642 & 117 \\
December  & 380 & 383 & 152 & 658 & 128 \\
\bottomrule
\end{tabular}
\end{table}


% ══════════════════════════════════════════════════════════════════════════════
\chapter{Full Counterfactual CO\texorpdfstring{\textsubscript{2}}{2} Scenario Results}
\label{app:counterfactual}

\Cref{tab:app_cf_wine,tab:app_cf_credit,tab:app_cf_higgs} give exact CO\textsubscript{2} values for all models under the eight carbon intensity scenarios introduced in \Cref{sec:results_carbon}. The \emph{Static (CC)} row is the hardware-corrected energy multiplied by CodeCarbon's built-in static intensity for Germany; the remaining seven rows use scenario intensities derived from the 2024 ElectricityMaps data. Wine and Credit values are in milligrams; HIGGS values are in grams.

\begin{table}[htbp]
\centering
\caption{Counterfactual CO\textsubscript{2} emissions (mg) for the Wine dataset under eight carbon intensity scenarios.}
\label{tab:app_cf_wine}
\begin{tabular}{@{} l r r r r r @{}}
\toprule
\textbf{Scenario} & \textbf{Log.\ Reg.} & \textbf{Rand.\ Forest} & \textbf{XGB (CPU)} & \textbf{XGB (GPU)} & \textbf{\gls{mlp}} \\
\midrule
Static (CC)    & 13.7 & 28.8 & 20.5 &  51.2 &  58.7 \\
Year mean      & 12.3 & 25.8 & 18.3 &  45.8 &  52.6 \\
Year best      &  4.1 &  8.7 &  6.2 &  15.5 &  17.7 \\
Year worst     & 23.7 & 49.8 & 35.3 &  88.5 & 101.4 \\
Winter mean    & 13.3 & 27.9 & 19.8 &  49.6 &  56.9 \\
Summer mean    & 10.9 & 23.0 & 16.3 &  40.8 &  46.8 \\
Summer noon    &  7.0 & 14.6 & 10.4 &  26.0 &  29.8 \\
Winter evening & 14.0 & 29.4 & 20.9 &  52.2 &  59.9 \\
\bottomrule
\end{tabular}
\end{table}

\begin{table}[htbp]
\centering
\caption{Counterfactual CO\textsubscript{2} emissions (mg) for the Credit dataset under eight carbon intensity scenarios.}
\label{tab:app_cf_credit}
\begin{tabular}{@{} l r r r r r @{}}
\toprule
\textbf{Scenario} & \textbf{Log.\ Reg.} & \textbf{Rand.\ Forest} & \textbf{XGB (CPU)} & \textbf{XGB (GPU)} & \textbf{\gls{mlp}} \\
\midrule
Static (CC)    &  60.5 &    831.6 & 43.3 &  51.7 &    633.0 \\
Year mean      &  54.2 &    744.2 & 38.8 &  46.3 &    566.5 \\
Year best      &  18.3 &    251.0 & 13.1 &  15.6 &    191.1 \\
Year worst     & 104.5 & 1{,}436.1 & 74.9 &  89.3 & 1{,}093.2 \\
Winter mean    &  58.6 &    805.6 & 42.0 &  50.1 &    613.2 \\
Summer mean    &  48.2 &    662.7 & 34.5 &  41.2 &    504.5 \\
Summer noon    &  30.7 &    422.2 & 22.0 &  26.3 &    321.4 \\
Winter evening &  61.7 &    848.1 & 44.2 &  52.7 &    645.6 \\
\bottomrule
\end{tabular}
\end{table}

\begin{table}[htbp]
\centering
\caption{Counterfactual CO\textsubscript{2} emissions (g) for the HIGGS dataset under eight carbon intensity scenarios. Random Forest was not run on the full 11M-row dataset.}
\label{tab:app_cf_higgs}
\begin{tabular}{@{} l r r r r @{}}
\toprule
\textbf{Scenario} & \textbf{Log.\ Reg.} & \textbf{XGB (CPU)} & \textbf{XGB (GPU)} & \textbf{\gls{mlp}} \\
\midrule
Static (CC)    &   5.148 &  14.650 &  2.104 & 145.892 \\
Year mean      &   4.607 &  13.110 &  1.883 & 130.558 \\
Year best      &   1.554 &   4.422 &  0.635 &  44.036 \\
Year worst     &   8.891 &  25.301 &  3.634 & 251.960 \\
Winter mean    &   4.987 &  14.192 &  2.038 & 141.329 \\
Summer mean    &   4.103 &  11.676 &  1.677 & 116.269 \\
Summer noon    &   2.614 &   7.438 &  1.068 &  74.075 \\
Winter evening &   5.251 &  14.941 &  2.146 & 148.792 \\
\bottomrule
\end{tabular}
\end{table}
